\documentclass[10pt,a4paper]{amsart}
\usepackage[margin=19mm]{geometry}
\usepackage{amsmath,amssymb,mathtools}
\usepackage[scr=rsfs,cal=euler]{mathalfa}
\usepackage{xfrac}
\usepackage[unicode,hidelinks,bookmarks=false]{hyperref}
\newcommand{\ie}{i.e.\ }
\newcommand{\cf}{cf.\ }
\newcommand{\resp}{respectively\ }
\newcommand{\Subsection}{\subsection}
\providecommand{\nobreakdash}{\nobreak\hbox{-}}
\setlength{\emergencystretch}{2em}
\multlinegap0pt
\usepackage{mathtools}
\usepackage{amssymb}
\usepackage{xcolor}
\usepackage[normalem]{ulem}
\usepackage{bm}
\newtheorem{lemma}{Lemma}
\newtheorem{theorem}{Theorem}
\def\F{{\mathbb F}}
\DeclareMathOperator{\Tr}{Tr}
\newcommand{\mc}{\mathcal}
\title{Factorizations of linearized polynomials and extremal curves in odd characteristic}
\author{Tetsushi Ito, Daichi Takeuchi, Takahiro Tsushima}
\date{Source: arXiv:2606.18577v1 (CC BY 4.0)}
\begin{document}
\maketitle

\noindent\emph{Source and licence.} Factorizations of linearized polynomials and extremal curves in odd characteristic. Version arXiv:2606.18577v1 (CC BY 4.0), distributed by arXiv under the Creative Commons Attribution 4.0 International licence, http://creativecommons.org/licenses/by/4.0/, as recorded in the arXiv licence metadata for this exact version and shown on the abstract page https://arxiv.org/abs/2606.18577v1. Full licence notice: \texttt{licenses/arxiv-2606.18577-CC-BY-4.0.txt}.\par\smallskip

\noindent\textit{Editorial scope.} Counted unit: the article's lemma key (source0.tex lines 1488--1503). The article's complete original proof of it is delivered with it (source0.tex lines 1504--1583, one \texttt{proof} environment of 80 lines). No mathematics is written, repaired or re-derived: the statement and the proof are the article's own lines, in the article's own notation, and no retained line is substituted or re-flowed.  Retained uncounted as the source-local context the proof needs: lines 496--521; lines 778--795; lines 983--1002; lines 1354--1366.  Nothing was omitted from any retained span, so the package carries no omission record.  No retained line contains a citation, so the wrapper carries no bibliography and no bibliography entry.  No retained line contains a figure environment, an image-inclusion command or drawing code, so no image is delivered and the package declares no asset dependency.  Editorial formatting only, in addition: an AMS \texttt{amsart} preamble that restates the article's own declarations and macros used by the retained lines, namely the environments \texttt{lemma}, \texttt{theorem} on separate counters, together with the standard packages those lines need.  The retained environments print in this document as Lemma 1, Theorem 1 in order of appearance, whereas the article prints the same items with the numbering of its own full document.  The complete native source of the article as distributed in the arXiv e-print for this exact version is bundled unchanged as \texttt{sources/203090/source0.tex}.
\begin{lemma}\label{lem: Fast F=R+Rast}
Assume that $p_0=2$. 
Let $F\in\mc R$ be an $\F_p$-linearized polynomial satisfying $F^\ast(1)=0$. Then the following hold. 
\begin{enumerate}
    \item There exists a unique $\F_p$-linearized polynomial $R(x)\in \F_q[x]$ of the form 
    \[
    R(x)=\sum_{i\geq1}a_ix^{p^i}\]
    such that 
\begin{equation*}
F^\ast \circ F=R+R^\ast. 
\end{equation*}
\item 
Let $R(x)$ be as in {\rm (i)}. Write 
\[
F(x)=\sum_{i=0}^eb_ix^{p^i}, 
\]
and define 
\[
\gamma_F\coloneqq\sum_{0\le i<j\le e}b_i^{p^{-i}}b_j^{p^{-j}}. 
\]
Then, in the ring $W_2(\F_q[x])$, we have 
\[
(F(x),0)\sim (0,x R(x)+\gamma_Fx^2). 
\]
\end{enumerate}
\end{lemma}

\begin{theorem}\label{t1}
Assume that $[\F_q:\F_p]$ is even. Then the following hold.  
\begin{itemize}
\item[{\rm (i)}]  Let $\rho=(W,a)$ be a pair as above. 
Then the curve $\overline{D}_\rho$ is $\F_q$-extremal if and only if 
\[\Tr_{q/p}(a^{-1}v^2)=0 \qquad\text{for all }v\in W.\]
Moreover, if these conditions are satisfied, then $\overline{D}_\rho$ is $\F_q$-maximal (resp. $\F_q$-minimal) if and only if 
    \[
    \left(\frac{a}{q}\right)\cdot \left(\frac{-1}{\sqrt{q}}\right)=-1 \qquad\text{(resp. }\left(\frac{a}{q}\right)\cdot \left(\frac{-1}{\sqrt{q}}\right)=1). 
    \]
\item[{\rm (ii)}] 

Let $R(x)\in\F_q[x]$ be an $\F_p$-linearized polynomial. Assume that  $\overline{C}_R$ is 
$\F_q$-extremal. Then there exists a pair $\rho=(W,a)$ satisfying 
\[R=R_\rho. \]
Consequently, we have $C_{R}=D_\rho$. 
\end{itemize}
\end{theorem}

\begin{theorem}\label{recipe for ex curve intro}
Assume that $[\F_q:\F_p]$ is even. Then the following statements hold: 
\begin{enumerate}
    \item The curve $\overline{D}_{\epsilon}$ is $\F_q$-extremal if and only if 
    \[
    Q_q(v)=\psi_q(tv)\qquad\text{for all } v\in W. 
    \]
    Moreover, if these conditions are satisfied, then 
    $\overline{D}_{\epsilon}$ is $\F_q$-maximal (resp.\ $\F_q$-minimal) if and only if 
    \[
    Q_q(t)=-(\sqrt{-1})^{[\F_q:\F_2]/2}\qquad\text{(resp.\ }Q_q(t)=(\sqrt{-1})^{[\F_q:\F_2]/2}\text{)}. 
    \]
    \item Let 
\[R(x)=\sum_{i=0}^e a_i x^{p^i} \in 
 \F_q[x],\]
 where $e\geq1$ and $a_e\neq0$, and assume that $\overline{C}_R$ is $\F_q$-extremal. Then there exists a pair $\epsilon=(W,t)$ as above and $a\in\F_q^\times$ such that $(x,y)\mapsto (ax,y)$ induces an isomorphism 
 \[ C_R\xrightarrow{\cong}D_{\epsilon}.
 \]
\end{enumerate}
\end{theorem}

\begin{lemma}\label{lemma: F,tF sim  xRx}
The following statements hold. 
\begin{enumerate}
    \item We use the notation in Lemma~\ref{lem: Fast F=R+Rast}. In the ring $W_2(\F_q[x])$, we have 
    \[
    \left(F_W(x),(tF_W(x))^2\right)\sim (0,xR(x)). 
    \]
    \item For every $x\in \F_q$, we have 
    \[
Q_q(F_W(x))\psi_q(tF_W(x))^{-1}=\psi_q(xR(x)). 
    \]
\end{enumerate}
\end{lemma}

\begin{lemma}\label{key}
The following statements hold. 
\begin{itemize}
\item[{\rm (i)}] 
Assume that $p$ is odd. 
Let $R(x)=ax$. Then $\overline{C}_{R}$ is 
$\F_q$-maximal (resp.\ $\F_q$-minimal) if and only if
\[
\left(\frac{a}{q}\right) \left(\frac{-1}{\sqrt{q}}\right)=-1 \quad (resp.\ \left(\frac{a}{q}\right) \left(\frac{-1}{\sqrt{q}}\right)=1). 
\]
\item[{\rm (ii)}] Assume that $p$ is even. 
Let $R_a(x)=x^p+ax$ for $a \in \F_q$. Then the family of curves 
$\{\overline{C}_{R_a}\}_{a \in \F_q}$ contains an $\F_q$-maximal curve. 
Furthermore, if $[\F_q:\F_p]>2$, then this family also contains an $\F_q$-minimal curve. 
\end{itemize}
\end{lemma}

\begin{proof}
(i) The assertion follows from 
Theorem~\ref{t1} applied 
to $(W,a)=(\{0\},a)$. 

\noindent
(ii) Set $W := \F_p$, and 
let $t \in \F_q$ satisfy $t^{\sqrt{q}} + t = 1$. Set $\epsilon:=(W,t)$. 

Since $W\subset\F_{\sqrt{q}}$, for every $v\in W$, we have
\[
Q_q(v)=\psi_{\sqrt{q}}(v)=\psi_q(tv). 
\]
By Theorem~\ref{recipe for ex curve intro}(i), the curve $\overline{D}_{\epsilon}$ is $\F_q$-extremal. 
Let $m\coloneqq [\F_q:\F_p]/2$. 
A direct computation shows that 
\begin{align*}
Q_q(t)
&=Q_{\sqrt{q}}(t^{\sqrt{q}}+t,t^{\sqrt{q}+1})\\
&=Q_{\sqrt{q}}(1,t+t^2)\\
&=\xi_2\!\left(m(1,0)+(0,\Tr_{\sqrt{q}/2}(t^2+t))\right)\\
&=(\sqrt{-1})^{m}\cdot \xi_2(0,t^{\sqrt{q}}+t)\\
&=(\sqrt{-1})^{m}\cdot \xi_2(0,1)\\
&=-(\sqrt{-1})^{m}.
\end{align*}
Hence, by Theorem~\ref{recipe for ex curve intro}(i),
$\overline{D}_{\epsilon}$ is $\F_q$-maximal.





 

Assume now that $m>1$. 
We show that there exists $u \in \F_q$ such that 
\begin{equation}\label{qqt}
Q_q(t + F_W(u)) = - Q_q(t).
\end{equation}

Since $Q_q(x + y) = Q_q(x)\, Q_q(y)\, \psi_q(xy)$ for $x,y \in \F_q$, we compute 
\begin{align*}
Q_q(t + F_W(u))Q_q(t)^{-1}
&= Q_q(F_W(u))\, \psi_q(tF_W(u)) \\
&=  \psi_q(uR_{\epsilon}(u)), 
\end{align*}
where the last equality follows from Lemma~\ref{lemma: F,tF sim  xRx}(ii). 

Assume that $\psi_q\bigl(u R_{\epsilon}(u)\bigr) = 1$ for all $u \in \F_q$. Then, for every $a\in\F_p$, 
\[
\psi(a\Tr_{q/p}(u R_{\epsilon}(u)))=
\psi_q\bigl(au R_{\epsilon}(u)\bigr) =\psi_q\bigl(a^{1/2}u R_{\epsilon}(a^{1/2}u)\bigr) = 1. 
\]
Hence, 
\[
\Tr_{q/p}\bigl(u R_{\epsilon}(u)\bigr) = 0
\quad \text{for all } u \in \F_q.
\]
By Artin--Schreier theory, it follows that 
\[
\# D_{\epsilon}(\F_q) =\# \left\{(u,v) \in \F_q^2 \mid 
v^p-v=u R_{\epsilon}(u)\right\}=pq. 
\]
On the other hand, since $\overline{D}_\epsilon$ is $\F_q$-maximal and has genus $p(p-1)/2$, we have 
\[
\# D_{\epsilon}(\F_q) = q + p(p-1)p^m. 
\]
Since $m > 1$, we have $q + p(p-1)p^m\neq pq$, a contradiction. 

Therefore, there exists $u \in \F_q$ such that 
\[
Q_q(t + F_W(u)) = - Q_q(t).
\]
Using \eqref{qqt}, for $v \in W=\ker F_W^\ast$, 
\begin{align*}
\psi_q((t+F_W(u))v)&=\psi_q(tv) \psi_q(F_W(u) v)\\
&=\psi_q(tv)\psi_q(u F_W^\ast(v))=\psi_q(tv)=Q_q(v). 
\end{align*}
By Theorem~\ref{recipe for ex curve intro}(i), the curve $\overline{D}_{(W,t+F_W(u))}$ is $\F_q$-minimal. 
\end{proof}

\end{document}
